\documentclass{article}
\setlength{\topmargin}{-0.25 in}
\setlength{\textwidth}{6.5in}
\setlength{\textheight}{8.5 in}
\setlength{\oddsidemargin}{0in}
\setlength{\evensidemargin}{0in}
\setlength{\parindent}{0.0in}
\setlength{\parskip}{7pt plus 1pt}

\usepackage{bcprules}
\usepackage{xcolor}				% Allow for color (annotations)

\def\tdm#1{\textcolor{blue}{\textsf{TDM: #1}}}
\def\ab#1{\textcolor{red}{\textsf{AB: #1}}}
\def\tf#1{\textcolor{purple}{\textsf{TF: #1}}}

\newcommand{\eval}[3]{{\tt #1} $\vdash$ {\tt #2} $\Downarrow$ {\tt #3}}
\newcommand{\solve}[2]{{\tt #1} $\models$ {\tt #2}}
\newcommand{\step}[5]{{\tt #1;#2} $\vdash$ {\tt #3} $\Longrightarrow$
  {\tt #4;#5}}
\newcommand{\tc}[3]{{\tt #1} $\vdash$ {\tt #2} : {\tt #3}}

\newcommand{\evalH}[4]{{\tt #1;#2} $\vdash$ {\tt #3} $\Downarrow$ {\tt #4}}
\newcommand{\solveH}[3]{{\tt #1;#2} $\models$ {\tt #3}}
\newcommand{\stepH}[7]{{\tt #1;#2;#3} $\vdash$ {\tt #4} $\Longrightarrow$
  {\tt #5;#6;#7}}
\newcommand{\tcH}[4]{{\tt #1;#2} $\vdash$ {\tt #3} : {\tt #4}}

\newcommand{\inferrule}[3]{\infrule[#1]{\mbox{#2}}{\mbox{#3}}}
\newcommand{\inferax}[2]{\infax[#1]{\mbox{#2}}}

\newcommand{\judgment}[1]{\begin{flushleft} \fbox{#1}\end{flushleft}}

\begin{document}

\title{A Formal Semantics for Babelsberg}
\author{Todd Millstein and Alan Borning and \ldots}
\date{\today}

\maketitle

\section{Introduction}

\emph{This draft is intended to evolve into a VPRI Memo.}

Babelsberg \cite{felgentreff-babelsberg-tr} is a family of object
constraint languages, with current instances being Babelsberg/R (a
Ruby extension), and Babelsberg/JS (a Javascript extension).  There
have been a number of confusing issues with respect to constraints and
object identity.  In an effort to understand these better, we are
constructing a formal operational semantics for Babelsberg.  The
results to date are promising, and we hope they will feed back into
the language design and implementation.

We approach this problem incrementally, first presenting a formal semantics
for a very simple constraint imperative language, and then building up
to a full object constraint language.  The languages are:

\begin{itemize}

\item Babelsberg/Int (the only datatype is integers)
\item Babelsberg/Records (Babelsberg with integers, booleans, and 
immutable records)
\item Babelsberg/MRecords (Babelsberg with integers, booleans, and 
mutable records)
\item Babelsberg/UID (Babelsberg with integers, booleans, and
  mutable records that live on the heap and so have an identity, can
  be aliased, etc.)
\item Babelsberg/Objects (Babelsberg with mutable objects, classes,
  methods, messages, and inheritance)
\end{itemize}

In each case, we first provide an informal discussion and examples,
and then the formal semantics.

We've found it helpful to approach the problem incrementally in
this fashion.  After the ideas have been refined, we plan to write a
paper for a conference or journal.  (This might only include a formal
semantics for Babelsberg/Objects, depending on whether it easily
stands on its own.)

\section{Constraints}

The semantics for each of these languages includes a step in which we
assert that some set of values for variables is a correct solution to
a collection of hard and soft constraints.  There is a constraint
solver that is a black box as far as the rest of the formal semantics
is concerned, and that handles all the issues regarding finding a
solution, dealing with conflicting soft constraints, and so forth.

We use the semantics for hard and soft constraints presented in
\cite{borning-lisp-symbolic-computation-1992}.  The JOT paper on
Babelsberg \cite{felgentreff-babelsberg-tr} that is currently under
review has a description of the relevant theory as well.  

The solver should find a single best solution --- if there are multiple
solutions, the solver is free to pick any one of them.  (Providing
answers rather than solutions, and backtracking among multiple
answers, as with constraint logic programming, is left for future work.)

The way we trade off conflicting soft constraints is defined by a
\emph{comparator}.  The two most relevant comparators are
locally-predicate-better (LPB) and weighted-sum-better (WSB)\@.
Locally-predicate-better only cares whether a constraint is satisfied
or not, not how far off the value is from the desired one.  Any
Pareto-optimal solution is OK\@.  For example, a solution that
satisfies one weak constraint A and violates three weak constraints B,
C, and D is OK, as long as there isn't a solution that satisfies both
A and some additional constraint; even if there is another solution
that satisfies B, C, and D but not A\@.  The DeltaBlue solver finds a
LPB solution.  Weighted-sum-better considers the weighted sum of the
errors of constraints with a given soft priority, and picks a solution
that minimizes the sum.  If there is more than one solution and there
are additional lower priority constraints, we then consider the
lower-priority ones to winnow down the possible solutions, priority by
priority.  For this comparator we need an error in satisfying a
constraint, which should be 0 iff the constraint is satisfied.
Cassowary finds WSB solutions.

Here are two examples.
\begin{verbatim}
required x+y=10
strong x=8
weak y=0
\end{verbatim}
The required constraint has an infinite number of solutions.  When we
winnow these down to solutions that satisfy the strong constraint,
there is only one left: $x=8$, $y=2$.  This is both a LPB and a WSB
solution.  The weak constraint has no impact on the solution in this
case.

Now consider:
\begin{verbatim}
required x+y=10
strong x>=5
weak y=20
\end{verbatim}

We'll only consider the WSB comparator this time, since it is more
suitable for use with inequalities.  (DeltaBlue does not handle inequalities.)
Again, the required constraint has an infinite number of solutions.
We winnow these down with the strong constraint to all $x \in
[5,\infty)$, $y$ such that $x+y=10$.   The weak constraint is
  unsatisfiable, so we minimize its error, resulting in the solution
  $x=5$, $y=5$.

Note that strict inequality constraints with metric comparators can be
problematic in the presence of soft constraints, since they can lead
to no solutions.  Consider:
\begin{verbatim}
required x>10
weak x=5
\end{verbatim}
This set of constraints has no solution --- for any potential solution
that satisfies \verb|x>10|, we can find another that better satisfies
\verb|x=5|.  For this reason, our examples usually use non-strict
inequalities. 

The ``required'' priority is special, in that those constraints must
be satisfied in any solution.  Both the semantics of hard and soft
constraints, and the Cassowary and DeltaBlue solvers, can handle
arbitrary numbers of soft priorities.  However, for simplicity in the
remainder of this note, we only use one, namely ``weak''.

We can also annotate variables used in constraints as
\emph{read-only}.  Intuitively, when choosing the best solutions to a
set of constraints with priorities, constraints should not be allowed
to affect the choice of values for their read-only variables, i.e.,
information can flow out of the read-only variables, but not into
them.  A \emph{one-way constraint} can be represented by annotating
all but one of the constrained variables as read-only.

In all of the languages, the constraint solver satisfies the
constraints at each time step, but does not itself need to know about
time.  Instead, time is built into the semantic rules; as far as the
solver is concerned, time is handled at each step by adding weak stay
constraints that variables keep their old values.  In other words, the
formal semantics handles time, while the solver handles solving the
constraints and dealing with hard and soft constraints.

\section{Babelsberg/Int}
\label{sec:babelsberg/int}

In this language, we only have one datatype: integers.  There is an
\emph{environment} that maps names to integer values.

Here are some examples.

\begin{verbatim}
x := 3;
x := 4;
always x>=10;
\end{verbatim}

After evaluating the first statement the environment has the binding
$x=3$.  After the second statement we hand the following constraints
to the solver to find a new value for $x$:
\begin{verbatim}
weak x=3 
required x=4
\end{verbatim}

The weak $x=3$ constraint is the stay constraint that $x$ retain its
previous value, while the required $x=4$ constraint comes from the
second assignment.  This has the solution $x=4$, which becomes the new
environment.  After the third statement, which adds the {\tt always}
constraint, we have the constraints:
\begin{verbatim}
weak x=4
required x>=10
\end{verbatim}

If we use a metric comparator such as weighted-sum-better (WSB) we get
the solution $x=10$, since this minimizes the error for the weak
constraint.  If we use LPB, then any $x \in [10,\infty)$ is a solution,
  and the system is free to select any of them.  (However, as noted
  above, typically we wouldn't use LPB if we have inequalities.)

Here is another example that involves using the same variable on the
right and left hand sides of an assignment statement.
\begin{verbatim}
x := 3;
x := x+2;
\end{verbatim}

After evaluating the first statement the environment has the binding
$x=3$.  After the second statement we hand the following constraints
to the solver to find a new value for $x$:
\begin{verbatim}
weak x=3 
required x=5
\end{verbatim}

Again, the weak $x=3$ constraint is the stay constraint that $x$
retain its previous value.  The required $x=5$ constraint comes from
the second assignment.  This is generated by evaluating {\tt x+2} in
the old environment (in which $x=3$), and then turning this into the
$x=5$ constraint.  After the second assignment, {\tt x} has the value
5, as expected.

An interesting aspect of this semantics, both as presented informally
above and in the formalism that follows, is that we no longer model
assignment as a constraint on variables at different times.  (In the
refinment model, as presented in the Babelsberg JOT paper
\cite{felgentreff-babelsberg-tr}, we described an assignment {\tt x :=
  x+2} as a constraint $x_{t+1} = x_t?  + 2$.)  But now the constraint
solver doesn't know anything about time.  Instead, we evaluate {\tt
  x+2} in the current environment to get 5.  Then we satisfy the
constraint {\tt x=5} in a new environment.

\ab{I \emph{think} the refinement model and the semantics here are
  equivalent.  Can we show this formally?  In particular, in the
  refinement model, when we represent an assignment as a constraint
  among variables at time $t$ on the right hand side, and a variable
  at $t+1$ on the left hand side, we annotate all the variables on the
  right-hand-side as read-only.  I think this is equivalent to just
  evaluating the expression on the right hand side, and setting up the
  constraint $v=c$ for a variable $v$ and constant $c$. Hmmm \ldots}

Depending on the solver, the program might introduce simultaneous
equations, general simplex problems, and so forth.  For example:
\begin{verbatim}
x := 0
y := 0
z := 0
always x+y+2*z = 10
always 2*x+y+z = 20
x := 100;
\end{verbatim}

After the final assignment we should have \verb|x=100, y=-270, z=90|.

\ab{We need to explain how this works! How does the set of constraints
  get constructed?  We've in effect got two different uses for +: one
  as the ordinary addition operator, the other as a constraint.
  Should we use the technique that the real Babelsbergs use, and run
  the {\tt always} constraint in constraint construction mode to get
  the constraints to hand off to the solver?}

As an example of unsatisfiable constraints, consider:
\begin{verbatim}
x := 5;
always x<=10;
x := x+15;
\end{verbatim}

After evaluating the first statement the environment includes the binding
$x=5$.  After evaluating the statement that generates the {\tt always}
constraint, we solve the constraints

\begin{verbatim}
weak x=5 
required x<=10
\end{verbatim}

This has the solution {\tt x=5}.  Then we evaluate the last
assignment, resulting the constraints
\begin{verbatim}
weak x=5 
required x<=10
required x=20
\end{verbatim}

These constraints are unsatisfiable.

\ab{What do we do now in the semantics? Should there be a ``bottom''
  value?  Presumably in the simpler languages we don't want to bother
  with exception handlers.}

This behavior is the same as that of the refinement model, in which
the program would be equivalent to these constraints (which are also
unsatisfiable): 
\begin{eqnarray*}
x_0 = 5                                \\
\forall t>0 \hspace*{2mm} x_t \leq 10  \\
\mbox{weak\ \ } x_1 = x_0?             \\
\mbox{weak\ \ } x_2 = x_1?             \\
x_2 = x_1? + 15
\end{eqnarray*}

It is \emph{different} from the behavior of the perturbation model (at
least one version of the perturbation model), in which we would assign
20 to x on the last statement, and then bump it down to 10 to satisfy
the {\tt always} constraint.

Interestingly, if we have read-only variables, and ``once''
constraints in addition to ``always'' constraints, we can simulate
assignment just with constraints --- we don't need an assignment statement.
need mutable records as such.  Consider the following example again:
\begin{verbatim}
x := 3;
x := 4;
\end{verbatim}

We can simulate this using {\tt once} constraints:

\begin{verbatim}
once x = 3;
once x = 4;
\end{verbatim}

If we have the same variable on the left and right hand sides of the
assignment, we need to use a temporary with a read-only annotation.
Thus:

\begin{verbatim}
x := 3;
x := x+2;
\end{verbatim}

becomes

\begin{verbatim}
once x = 3;
once temp = x?;
once x = temp?+2;
\end{verbatim}

\subsection{Formalism}
\label{intformal}

\subsubsection{Syntax}

A program is a statement defined by the
following syntax:

\begin{tt}
\begin{tabular}{lrrl}
{\rm Statement} & s & ::= & x := e $\mid$ [weak] always C $\mid$ s;s \\
{\rm Expression} & e & ::= & b $\mid$ n $\mid$ x $\mid$ e op e \\
{\rm Constraint} & C & ::= & e \\
{\rm Boolean} & b & ::= & true $\mid$ false \\
{\rm Number} & n & ::= & {\rm numeric constants} \\
{\rm Variable} & x & ::= & {\rm variable names} \\
{\rm Operator} & op & ::= & + $\mid$ - $\mid$ >= $\mid$ $\cdots$ \\\\

{\rm Value} & v & ::= & b $\mid$ n
\end{tabular}
\end{tt}

\subsubsection{Semantics}

The semantics is defined by several judgments, defined below.  These judgments
depends on the notion of an {\em environment}, which is a partial
function from variables to values.  Metavariable {\tt E} ranges over
environments.  We also assume that each operator {\tt op} comes with
an associated semantic function, denoted {\tt [[op]]}.

\judgment{\eval{E}{e}{v}}

``Expression {\tt e} evaluates to value {\tt v} in the context of environment {\tt
  E}.''

\inferax{E-Bool}{\eval{E}{b}{b}}

\inferax{E-Num}{\eval{E}{n}{n}}

\inferrule{E-Var}{{\tt E}({\tt x}) = {\tt v}}{\eval{E}{x}{v}}

\inferrule{E-Op}{\eval{E}{e$_1$}{v$_1$} \andalso
\eval{E}{e$_2$}{v$_2$} \andalso
{\tt v$_1$ [[op]] v$_2$} = {\tt v}}
{\eval{E}{e$_1$ op e$_2$}{v}}

\judgment{\solve{E}{C}}

``Environment {\tt E} is a {\em model} of constraint {\tt C}.''

This judgment represents a call to the constraint solver, which we
treat as a black box.  In particular, the proposition 
\solve{E}{C} denotes that {\tt E} is a solution to the constraints
{\tt C} (and further one that is optimal according to the solver's
semantics, as discussed earlier).  We assume that the logic of the
solver supports solving conjunctions of arithmetic constraints, with
each arithmetic constraint optionally labeled with a priority.
\tdm{Please check, Alan!}
\ab{Reworded as above - was: ``We assume that the logic of the
solver includes propositional logic with arithmetic constraints and
that it supports the distinction between strong and weak
constraints.'' Sound OK?}

\judgment{\step{E}{C}{s}{E$'$}{C$'$}}

``Statement {\tt s} transforms environment {\tt E} and global
constraint {\tt C} to environment {\tt E$'$} and global constraint
{\tt C$'$}.''

A ``state'' of execution includes a concrete portion, represented by
the environment, and a symbolic portion, represented by the
constraint.  The environment is standard, while the constraint is not
part of the state of a computation in most languages.  Intuitively, the environment comes from constraint solving during the
evaluation of the immediately preceding statement, and the 
constraint records the ``always'' constraints that have been
declared so far during execution. 

In the rules below we use the notation stay({\tt E}) to denote a
conjunction of constraints of the form {\tt weak x=v} for
each {\tt x} in the domain of {\tt E}, where {\tt E}({\tt x}) = {\tt v}.

\inferrule{S-Asgn}{\eval{E}{e}{v} \andalso
\solve{E$'$}{(C $\wedge$ {\rm stay({\tt E})} $\wedge$ x=v)}}
{\step{E}{C}{x := e}{E$'$}{C}}


\inferrule{S-Always}{
{\tt C$'$} = {\tt (C $\wedge$ [weak] C$_0$)} \andalso
\solve{E$'$}{(C$'$ $\wedge$ {\rm stay({\tt E})})}}
{\step{E}{C}{[weak] always C$_0$}{E$'$}{C$'$}}


\inferrule{S-Seq}{\step{E}{C}{s}{E$_0$}{C$_0$} \andalso
\step{E$_0$}{C$_0$}{s$'$}{E$'$}{C$'$}}
{\step{E}{C}{s;s$'$}{E$'$}{C$'$}}

\section{Babelsberg/Records}

For this next language, we augment the Babelsberg/Int language with
booleans and immutable records.

Records are represented as lists of name/value pairs in curly braces:
\begin{verbatim}
  {x:5, y:10}
\end{verbatim}

The fields can be either integers or other records, which of course
lets us make lists, trees, etc.  For example here is the list [1,2,3]:

\begin{verbatim}
  {head: 1, tail: {head: 2, tail: {head: 3, tail: {}}}}
\end{verbatim}

There is no notion of object identity Babelsberg/Records --- we can
test whether two records are equal, but whether or not they are
identical is an implementation issue and not part of the semantics.

The syntax is extended to include record constructors and field
access.  Here are examples of expressions, assignment statements, and
constraints involving records:

\begin{verbatim}
 p := {x:2, y:5};   /* assign a record to p */
 a := p.x;          /* access a field of a record and assign it to a variable */
 q := p;            /* q is now a copy of p (or maybe it's shared; we can't tell
                       the difference) */
 always p.x = 100;  /* a constraint on a record field */
 always r = p;      /* an equality constraint between two records */
\end{verbatim}

There is no static type checking - checking is all done dynamically.
Our model includes structural compatibility checks on constraints.
These allow us to desugar all of the record constraints (recursively)
into constraints on the individual components.  Thus, the record
constructs are all handled by the semantic rules, and the constraint
solver does not need to know anything about records.  The alternative
of having the solver itself know about records was rejected, since it
is more complex and seems to lead to more indeterminism, as well as 
some odd cases (including
requiring that the constraint solver spontaneously generate new
records).  For completeness, we outline this rejected alternative in 
Appendix \ref{sec:structural-compatability-check-alternatives}.

The structural compatibility checks are assertions that are checked
before doing desugaring records that check whether a variable holds a
record, and whether the record has the necessary fields.  These
assertions are checked, but unlike normal constraints
the system doesn't change anything to
enforce them --- if one is violated it's just an error.  Instead, the
programmer must ensure that a record with the expected fields is
assigned to a variable used in record constraints, just as a
programmer would need to ensure that a record with the expected fields
was assigned to a record-valued variable in a standard language.

Here are a few examples.

\begin{verbatim}
 p := {x:2, y:5};
 always p.x = 100;
\end{verbatim}
After executing the two statements, \verb|p| is \verb|{x:100, y:5}|.
Here is how the system determines this result.  The structural
compatibility check for the \verb|always| constraint is that \verb|p|
has an \verb|x| field.  We desugar the record constraint, so that
after evaluating the \verb|always p.x = 100| statement the constraints
are:
\begin{verbatim}
 weak px=2
 weak py=5
 required px = 100
\end{verbatim}
These constraints have the solution \verb|px=100, py=5|, which is then
reassembled into \verb|p = {x:100, y:5}|.

Here are a few examples that fail the structural compatibility checks.

\begin{verbatim}
 p := {x:2, y:5};
 always q = p;
\end{verbatim}
There is a check before the \verb|always| constraint that \verb|q|
have \verb|x| and \verb|y| fields, which fails.  To fix this, we could
initialize \verb|q| appropriately:
\begin{verbatim}
 p := {x:2, y:5};
 q := p;
 always q = p;
\end{verbatim}

(Experience with Babelsberg/R indicates that requiring this
initialization is annoying.  We can avoid the problem, as is described
in Section \ref{sec:record-convenience-features}.  However, doing so
clutters up the semantics a bit, and since our goal with
Bablesberg/Records is to illuminate the semantics rather than to
design a practical language, we don't include this in our formal
semantics.)

Similarly, we can't simply assert a constraint on a record \verb|p|:
\begin{verbatim}
 always p = {x:2, y:5};
\end{verbatim}

Rather, we need to initialize \verb|p| first:
\begin{verbatim}
 p := {x:2, y:5};
 always p = {x:2, y:5};
\end{verbatim}

The following program passes the structural compatibility checks, but
fails with an unsatisfiable constraint error, since we are requiring
that x be both 100 and 2:
\begin{verbatim}  
 p := {x:0, y:0};
 always p.x = 100;
 p := {x:2, y:5};
\end{verbatim}

However, the following program is OK:
\begin{verbatim} 
 p := {x:0};
 always weak p = {x:3};
 always weak p = {x:4}
\end{verbatim}

This desugars to the constraints \verb|weak px=3| and
\verb|weak px=4|, plus \verb|weak px=0| representing the stay
constraint.  Using the LPB comparator, there are three possible
answers: \verb|p={x:0}|, \verb|p={x:3}| or \verb|p={x:4}|.  A WSB
comparator would allow x to take on any value in $[0,4]$, although
Cassowary would produce either 1, 3 or 4 (but not something in
between), i.e., the same solutions as LPB.  (If we made the
\verb|always| constraints slightly higher priority than the stay, the
solutions would be either \verb|p={x:3}| or \verb|p={x:4}|.)

Here's another example that fails a structural compatibility check:
\begin{verbatim}
 a := {x:0};
 b := {y:0};
 always weak a.x = 3
 always weak b.y = 10
 always weak a=b
\end{verbatim}
This is the case even though there actually are records that would
satisfy all of the weak constraints.  Thus, if the initial assignments
had been different:
\begin{verbatim}
 a := {x:0, y:0};
 b := {x:0, y:0};
 always weak a.x = 3
 always weak b.y = 10
 always weak a=b
\end{verbatim}
the program works.

\subsection{Formalism}
\label{recformal}

\subsubsection{Syntax}

The syntax from Section~\ref{intformal} is augmented now to support
records and record operations:

\begin{tt}
\begin{tabular}{lrrl}
{\rm Expression} & e & ::= & $\cdots$ $\mid$
\{l$_1$:e$_1$,$\ldots$,l$_n$:e$_n$\} $\mid$ e.l \\
{\rm Label} & l & ::= & {\rm record label names} \\\\

{\rm Value} & v & ::= & $\cdots$ $\mid$
\{l$_1$:v$_1$,$\ldots$,l$_n$:v$_n$\}
\end{tabular}
\end{tt}

\subsubsection{Semantics}

\judgment{\eval{E}{e}{v}}

Expression evaluation is updated in the standard way to support records.

\inferrule{E-Rec}{\eval{E}{e$_1$}{v$_1$} $\cdots$
  \eval{E}{e$_n$}{v$_n$}}
{\eval{E}{\{l$_1$:e$_1$,$\ldots$,l$_n$:e$_n$\}}{\{l$_1$:v$_1$,$\ldots$,l$_n$:v$_n$\}}}

\inferrule{E-Field}{\eval{E}{e}{\{l$_1$:v$_1$,$\ldots$,l$_n$:v$_n$\}}
    \andalso 1 $\leq$ i $\leq$ n}
{\eval{E}{e.l$_i$}{v$_i$}}

\judgment{\solve{E}{C}}

This judgment represents the solver, as before.  However, now
constraints can refer to records and record operations, and we do not want to assume that the solver
directly supports records.  Therefore, all constraints on records get
``exploded'' into constraints on the individual components of the
records.  \tdm{Should we define this transformation more formally?}
For example, the constraint {\tt p = \{x:2, y:5\}} gets translated to
the constraint {\tt px = 2 $\wedge$ py=5}, and similarly the
constraint {\tt p.x = 100} becomes {\tt px = 100}.

\judgment{\tc{E}{e}{T}}

``Expression {\tt e} has type {\tt T} in the context of environment
{\tt E}.''

We use a nonstandard form of typechecking to perform the structural
compatibility check.  It is nonstandard because it happens
dynamically on an expression, in the context of the current run-time
environment.  Otherwise the rules are straightforward.  We use the
following syntax for types:

\begin{tt}
\begin{tabular}{lrrl}
{\rm Type} & T & ::= & Num $\mid$ Bool $\mid$ 
\{l$_1$:T$_1$,$\ldots$,l$_n$:T$_n$\} 
\end{tabular}
\end{tt}

We assume that each operator {\tt op} comes with
an associated typing function, denoted typeof$_{\tt op}$, which is a
partial function from a pair of types (representing the types of the
operands) to a type (representing the type of the result).

\inferax{T-Bool}{\tc{E}{b}{Bool}}

\inferax{T-Num}{\tc{E}{n}{Num}}

\inferrule{T-Var}{{\tt E}({\tt x}) = {\tt v} \andalso \tc{E}{v}{T}}{\tc{E}{x}{T}}

\inferrule{T-Op}{\tc{E}{e$_1$}{T$_1$} \andalso
\tc{E}{e$_2$}{T$_2$} \andalso
typeof$_{\tt op}$({\tt T$_1$}, {\tt T$_2$}) = {\tt T}}
{\tc{E}{e$_1$ op e$_2$}{T}}

Note: The intention is that typeof$_{\tt =}$ will only allow two
records to be compared for equality if they have the same type.  This
prevents nondeterminism that would otherwise occur in a constraint
that equates two record variables {\tt p} and {\tt q}.

\inferrule{T-Rec}{\tc{E}{e$_1$}{T$_1$} $\cdots$
  \tc{E}{e$_n$}{T$_n$}}
{\tc{E}{\{l$_1$:e$_1$,$\ldots$,l$_n$:e$_n$\}}{\{l$_1$:T$_1$,$\ldots$,l$_n$:T$_n$\}}}

\inferrule{E-Field}{\tc{E}{e}{\{l$_1$:T$_1$,$\ldots$,l$_n$:T$_n$\}}
    \andalso 1 $\leq$ i $\leq$ n}
{\tc{E}{e.l$_i$}{T$_i$}}

The above rule is a key one, ensuring that a field of a record can only be
referred to if it already exists in that record.

\judgment{\step{E}{C}{s}{E$'$}{C$'$}}

The semantics for executing statements is identical to what we had before,
except that now we ``typecheck'' an {\tt always} constraint to ensure
that it meets our structural compatibility requirements.  Therefore,
we simply keep the rules {\sc S-Asgn} and {\sc S-Seq} as is but replace the rule {\sc S-Always} with the following one,
which adds the compatibility check:


\inferrule{S-Always}{\tc{E}{C$_0$}{Bool} \andalso
{\tt C$'$} = {\tt (C $\wedge$ [weak] C$_0$)} \andalso
\solve{E$'$}{(C$'$ $\wedge$ {\rm stay({\tt E})})}}
{\step{E}{C}{[weak] always C$_0$}{E$'$}{C$'$}}

The compatibility check also ensures that a
constraint is in fact a boolean-valued expression.

\subsection{Convenience Features}
\label{sec:record-convenience-features}

As illustrated by some of the above examples, we need to initialize a
record variable with a suitable value before using it in an
\verb|always| constraint, even if the \verb|always| constraint would
define its entire structure.  As mentioned above, in practice this is
annoying, and so in the actual Babelsberg/R language we allow an
\verb|always| constraint to define the structure of an uninitialized
variable.  If we supported this in Babelsberg/Records as well, this
program would be legal:
\begin{verbatim}
 p := {x:2, y:5};
 always q = p;
\end{verbatim}

as would this one:
\begin{verbatim}
 always p = {x:2, y:5};
\end{verbatim}

\section{Babelsberg/MRecords}

Our next language is Babelsberg/MRecords, which involves extending
Babelsberg/Records to allow mutable records.  Syntactically, we simply
allow field acccesses as l-values, e.g.
\begin{verbatim}
 p := {x:0, y:0};
 p.x := 100;
\end{verbatim}

After the second assignment, this desugars to the constraints
\verb|weak px=0|, \verb|weak py=0|, \verb|px=100|, which after putting
the record back together has the solution \verb|p={x:100, y:0}|.  This
doesn't add any particular complications to the semantics.

There is still no notion of object identity.  Consider:
\begin{verbatim}
 p := {x:0, y:0};
 q := p;
 p.x := 100;
\end{verbatim}
Now \verb|p={x:100, y:0}|, but \verb|q={x:0, y:0}|.

\subsection{Formalism}
\label{mrecformal}

\subsubsection{Syntax}

The syntax from Section~\ref{recformal} is augmented now to support
assignment of record components:

\begin{tt}
\begin{tabular}{lrrl}
{\rm Statement} & s & ::= & $\cdots$ $\mid$ x.l := e
\end{tabular}
\end{tt}

\subsubsection{Semantics}

The semantics from Section~\ref{recformal} is similarly augmented now
to support assignment of record components.  That simply requires the
addition of one new rule to the semantics of statement evaluation:

\judgment{\step{E}{C}{s}{E$'$}{C$'$}}

\infrule[S-AsgnField]{\mbox{\tc{E}{x}{\{l$_1$:T$_1$,$\ldots$,l$_n$:T$_n$\}}
    \andalso 1 $\leq$ i $\leq$ n \andalso \eval{E}{e}{v}} \\
\mbox{\solve{E$'$}{(C $\wedge$ {\rm stay({\tt E})} $\wedge$ x.l$_i$=v)}}}
{\mbox{\step{E}{C}{x.l$_i$ := e}{E$'$}{C}}}

Note that we perform a structural compatibility check to ensure that
the field being mutated already exists in the record denoted by {\tt x}.

\subsection{Simulating Mutable Records}

If we have read-only variables, and ``once'' constraints in addition
to ``always'' constraints, we can simulate assigning to the field of a
record just with constraints --- we don't need mutable records as
such.  (This is similar to the simulation of assignment discussed for
Babelsberg/Int in Section \ref{sec:babelsberg/int}.)
Consider:

\begin{verbatim}
 p := {x:0, y:0};
 once p.x = 100;
\end{verbatim}

The \verb|once| constraint takes the place of an assignment to
\verb|p.x|.  If the right hand side of the assignment is a variable or
an expression, we need to annotate the variable (or all the variables
in the expression) as read-only to achieve the directionality of the
assignment. For example:

\begin{verbatim}
 p := {x:0, y:0};
 once p.x = z?;
\end{verbatim}

(Without the read-only annotation on \verb|z|, the solver might change
\verb|z| rather than \verb|p.x| to satisfy the constraint.)

If the assignment involves the same field of the record, we need to
use a temporary.  Thus, this program:

\begin{verbatim}
 p := {x:0, y:0};
 p.x := p.x+1;
\end{verbatim}

becomes

\begin{verbatim}
 p := {x:0, y:0};
 once temp = p.x?
 once p.x = temp?
\end{verbatim}

However, unlike Babelsberg/Int, we can't do away with assignment
altogether: it is needed to provide the initial structure for
variables that refer to records.  So this version would fail the
structural compatibility check:

\begin{verbatim}
 once p = {x:0, y:0};
 once temp = p.x?
 once p.x = temp?
\end{verbatim}

\ab{I want to move all of the discussion about once constraints and
  assignment to an appendix.  Also we could solve the structural
  compatibility check problem by using the convenience feature
  described before.}

\section{Babelsberg/UID}

Our next step is to replace records that are simply values with
mutable records that live on the heap and so have an identity, can be
aliased, and so forth.  (For simplicity, we drop records that are
simply values from Babelsberg/UID, although one could have a language
with both kinds.)  To emphasize that we now allocate records on the
heap, we use the keyword \verb|new| when creating one:

\begin{verbatim}
 p := new {x:0, y:0};
\end{verbatim}

Now variables in the environment that refer to records hold pointers,
which in turn reference records in the heap.  (This only applies to
variables that have a record type; integers and booleans don't live in
the heap.)  The following example illustrates the difference from
Babelsberg/MRecords:

\ab{Todd: please check wording above - I'm assuming integers and
  booleans don't live in the heap.}

\begin{verbatim}
 p := new {x:0, y:0};
 q := p;
 p.x := 100;
\end{verbatim}
Unlike in the previous example, now \verb|p| and \verb|q| are aliased,
so after the program runs we have \verb|p=q={x:100,y:0}|, with
\verb|p| and \verb|q| identical (not just equal).

Also unlike the previous examples, the weak stays on variables that
refer to records are identity constraints rather than equality
constraints: such a variable continues to refer to the same object
unless reassigned.  

\subsection{Identity Constraints}

\ab{I think the formal semantics doesn't account for identity
  constraints --- we need to update it accordingly.}


To make this work, 

Perhaps we can have identity constraints, in addition to equality constraints:
\begin{verbatim}
 p := new {x:0, y:0};
 always p==q;  /* where == is object identity */
 q := new {a:10};
\end{verbatim}
Now both \verb|p| and \verb|q| refer to the same record
\verb|{a:10}|.  (If the constraint had been an equality constraint
rather than an identity constraint, the assignment to \verb|q| would
have failed the structural compatibility check, which would have
expected it to be a record with \verb|x| and \verb|y| fields.)

\tf{shouldn't assignment be equivalent to a directed ``once'' identity
  constraint rather than an equality constraint now?} 
\ab{If we continue with the semantics we used for Babelsberg/Int, we
  aren't actually modeling assignment as a once constraint ---
  instead, object ID on the right hand side of the assignment is a
  constant, and just assigned to the variable on the left.}

To keep things simple, it may be useful to have a convention that
identity constraints are either just checks when they are first
created, or at least one side of the identity is a fresh variable.
Otherwise we may
have unexpectedly unsatisfiable identity constraints.  Consider:

\begin{verbatim}
 p := new {x:0, y:0};
 q := new {z:0};
 always p.x = 10;
 always q.z = 100;
 always p==q; 
\end{verbatim}
Up to the final identity constraint, this should be fine.  But the
final constraint is unsatisfiable (unless the solver creates a new
object, which isn't allowed with the current semantics).

We might also have priorities on identity constraints:
\begin{verbatim}
 p := new {x:0, y:0};
 q := new {z:0};
 always p.x = 10;
 always q.z = 100;
 always p==r; 
 always weak q==r; 
\end{verbatim}
Here, the weak identity constraint \verb|always weak q==r| can't be
satisfied.

\subsection{Formalism}

Since this version is significantly different than the prior formalisms, we present
it in its entirety rather than as a delta from those ones.

\subsubsection{Syntax}

The syntax is modified in a few ways, all standard.  First, records are not
expressions and there are also no record values.  Instead records are
created via  a {\tt new} statement, which allocates the record on the heap
and returns a {\em reference} to it.  We introduce
a metavariable {\tt r} ranging over mutable references, which are a new
kind of value in the language.  It is convenient to consider
references to be expressions, but we assume they are never used in
source programs.  \tdm{Not sure if we need that assumption.}

\begin{tt}
\begin{tabular}{lrrl}
{\rm Statement} & s & ::= & x := e $\mid$ x.l := e $\mid$ x := new \{l$_1$:e$_1$,$\ldots$,l$_n$:e$_n$\}
$\mid$ [weak] always C $\mid$ s;s \\
{\rm Expression} & e & ::= & b $\mid$ n $\mid$ x $\mid$ e op e $\mid$
e.l $\mid$ r \\
{\rm Constraint} & C & ::= & e \\
{\rm Boolean} & b & ::= & true $\mid$ false \\
{\rm Number} & n & ::= & {\rm numeric constants} \\
{\rm Variable} & x & ::= & {\rm variable names} \\
{\rm Operator} & op & ::= & + $\mid$ - $\mid$ >= $\mid$ $\cdots$ \\
{\rm Label} & l & ::= & {\rm record label names} \\\\

{\rm Value} & v & ::= & b $\mid$ n $\mid$ r \\
{\rm Reference} & r & ::= & {\rm mutable references}
\end{tabular}
\end{tt}

\subsubsection{Semantics}

In addition to an environment, the semantics now requires a {\em
  heap}, which is a partial function from mutable references to record
``objects'' of the form {\tt \{l$_1$:v$_1$,$\ldots$,l$_n$:v$_n$\}}.
Heaps are ranged over by metavariable {\tt H}.
Most of the updates to the original rules are standard and straightforward.

\judgment{\evalH{E}{H}{e}{v}}

``Expression {\tt e} evaluates to value {\tt v} in the context of environment {\tt
  E} and heap {\tt H}.''

\inferax{E-Bool}{\evalH{E}{H}{b}{b}}

\inferax{E-Num}{\evalH{E}{H}{n}{n}}

\inferrule{E-Var}{{\tt E}({\tt x}) = {\tt v}}{\evalH{E}{H}{x}{v}}

\inferrule{E-Op}{\evalH{E}{H}{e$_1$}{v$_1$} \andalso
\evalH{E}{H}{e$_2$}{v$_2$} \andalso
{\tt v$_1$ [[op]] v$_2$} = {\tt v}}
{\evalH{E}{H}{e$_1$ op e$_2$}{v}}

\inferrule{E-Field}{\evalH{E}{H}{e}{r} \andalso {\tt H}({\tt r}) =
  {\tt \{l$_1$:v$_1$,$\ldots$,l$_n$:v$_n$\}} 
    \andalso 1 $\leq$ i $\leq$ n}
{\evalH{E}{H}{e.l$_i$}{v$_i$}}

\inferax{E-Ref}{\evalH{E}{H}{r}{r}}

\judgment{\solveH{E}{H}{C}}

``Environment {\tt E} and heap {\tt H} are a model of constraint {\tt C}.''

As before, this judgment represents a call to the constraint solver.
We again ``explode'' constraints on records into constraints on the
individual components, but in a slightly different way than before due
to their heap allocation.  
For example, suppose variable {\tt p} maps to reference {\tt r} in the
environment.  Then the constraint
 {\tt p = \{x:2, y:5\}} gets translated to
the constraint {\tt rx = 2 $\wedge$ ry=5}, and similarly the
constraint {\tt p.x = 100} becomes {\tt rx = 100}.  By using the
reference {\tt r} rather than the variable name {\tt p}, we ensure
that any modifications to the record stored in {\tt r} will also be
``seen'' through all aliases of {\tt p}.  \tdm{More formal definition
  of our approach?}


\judgment{\tcH{E}{H}{e}{T}}

``Expression {\tt e} has type {\tt T} in the context of environment
{\tt E} and heap {\tt H}.''


\inferax{T-Bool}{\tcH{E}{H}{b}{Bool}}

\inferax{T-Num}{\tcH{E}{H}{n}{Num}}

\inferrule{T-Var}{{\tt E}({\tt x}) = {\tt v} \andalso \tcH{E}{H}{v}{T}}{\tcH{E}{H}{x}{T}}

\inferrule{T-Op}{\tcH{E}{H}{e$_1$}{T$_1$} \andalso
\tcH{E}{H}{e$_2$}{T$_2$} \andalso
typeof$_{\tt op}$({\tt T$_1$}, {\tt T$_2$}) = {\tt T}}
{\tcH{E}{H}{e$_1$ op e$_2$}{T}}


\inferrule{E-Field}{\tcH{E}{H}{e}{\{l$_1$:T$_1$,$\ldots$,l$_n$:T$_n$\}}
    \andalso 1 $\leq$ i $\leq$ n}
{\tcH{E}{H}{e.l$_i$}{T$_i$}}

\inferrule{T-Ref}{{\tt H}({\tt r}) = {\tt
    \{l$_1$:v$_1$,$\ldots$,l$_n$:v$_n$\}} \andalso \tcH{E}{H}{v$_1$}{T$_1$} $\cdots$
  \tcH{E}{H}{v$_n$}{T$_n$}}
{\tcH{E}{H}{r}{\{l$_1$:T$_1$,$\ldots$,l$_n$:T$_n$\}}}

Note that the above rule does not work as intended in the presence of
cyclic records, which are now possible.  In that case, we will simply
not be able to find a finite derivation and so will raise an error
even when the given record satisfies our intended structural
compatibility check.  This should be fixable by keeping track of
references that have already been ``visited'' in order to know when to
stop, and by adding a kind of recursive record type to represent the
structure of cyclic records.


\judgment{\stepH{E}{H}{C}{s}{E$'$}{H$'$}{C$'$}}

``Statement {\tt s} transforms environment {\tt E}, heap {\tt H}, and global
constraint {\tt C} to environment {\tt E$'$}, heap {\tt H'}, and global constraint
{\tt C$'$}.''

In the rules below, stay({\tt E}) denotes, as earlier,  a
conjunction of constraints of the form {\tt weak x=v} for
each {\tt x} in the domain of {\tt E}, where {\tt E}({\tt x}) = {\tt v}.
The notation stay({\tt H}) is analogous except that these record
constraints are ``exploded'' as discussed earlier.  For example, if
{\tt H}({\tt r}) = {\tt \{x:2, y:5\}}, then the constraints
{\tt weak rx = 2} and {\tt weak ry=5} are included in stay({\tt H}).
Finally, dom({\tt H}) denotes the domain of {\tt H}.

\inferrule{S-Asgn}{\evalH{E}{H}{e}{v} \andalso
\solveH{E$'$}{H$'$}{(C $\wedge$ {\rm stay({\tt E})} $\wedge$ {\rm stay({\tt H})} $\wedge$ x=v)}}
{\stepH{E}{H}{C}{x := e}{E$'$}{H$'$}{C}}

\infrule[S-AsgnField]{\mbox{\tcH{E}{H}{x}{\{l$_1$:T$_1$,$\ldots$,l$_n$:T$_n$\}}
    \andalso 1 $\leq$ i $\leq$ n \andalso \evalH{E}{H}{e}{v}} \\ 
\mbox{\solve{E$'$}{(C $\wedge$ {\rm stay({\tt E})} $\wedge$ {\rm stay({\tt H})} $\wedge$ x.l$_i$=v)}}}
{\mbox{\stepH{E}{H}{C}{x.l$_i$ := e}{E$'$}{H$'$}{C}}}


\infrule[S-AsgnNew]{\mbox{\evalH{E}{H}{e$_1$}{v$_1$} $\cdots$ \evalH{E}{H}{e$_n$}{v$_n$}
    \andalso r $\not\in$ dom({\tt H})} \\ 
\mbox{\solveH{E$'$}{H$'$}{(C $\wedge$ {\rm stay({\tt E})} $\wedge$ {\rm
    stay({\tt H})} $\wedge$ x=r $\wedge$ r = \{l$_1$:v$_1$,$\ldots$,l$_n$:v$_n$\})}}}
{\mbox{\stepH{E}{H}{C}{x := new \{l$_1$:e$_1$,$\ldots$,l$_n$:e$_n$\}}{E$'$}{H$'$}{C}}}

Recall that the constraint {\tt r =
  \{l$_1$:e$_1$,$\ldots$,l$_n$:e$_n$\}} in the above rule is shorthand for $n$
constraints of the form {\tt rl$_i$ = v$_i$}.

\inferrule{S-Always}{\tcH{E}{H}{C$_0$}{Bool} \andalso
{\tt C$'$} = {\tt (C $\wedge$ [weak] C$_0$)} \andalso
\solveH{E$'$}{H$'$}{(C$'$ $\wedge$ {\rm stay({\tt E})} $\wedge$ {\rm stay({\tt H})})}}
{\stepH{E}{H}{C}{[weak] always C$_0$}{E$'$}{H$'$}{C$'$}}

\inferrule{S-Seq}{\stepH{E}{H}{C}{s}{E$_0$}{H$_0$}{C$_0$} \andalso
\stepH{E$_0$}{H$_0$}{C$_0$}{s$'$}{E$'$}{H$'$}{C$'$}}
{\stepH{E}{H}{C}{s;s$'$}{E$'$}{H$'$}{C$'$}}

\section{Babelsberg/Objects}

In this final language, we add classes, inheritance, methods, and
messages.  This allows us to write constraints on the results of
message sends, as in the actual languages Babelsberg/R and
Babelsberg/JS\@.  For example constraints on the center of a
rectangle r:
\begin{verbatim}
always r.center.x = 100;
always r.center.y = 50;
\end{verbatim}
apply to the result of sending the message \verb|center| to \verb|r|,
and \verb|center| may be computed rather than being stored.  Even the
\verb|x| message to a point might compute the x value, if the point is
stored in polar coordinates.  

\emph{In the implementation, methods invoked in a constraint are
  evaluated in constraint construction mode, so that the system builds
  up a network of primitive constraints that represents the
  user-defined method.  The semantics should presumably mirror this
  process.}
\tf{If we want to mirror the way it works in the current implementation
with cooperating solvers, then we need to extend the semantics of how
a variable is evaluated to a value. In constraint construction mode,
this is a function over the variable \emph{and} the active solver,
that evaluates to either the normal value or, if the solver supports
the variables' type, to a solver-specific representation of the
value. Messages sent to a variable in constraint construction mode
thus may be called on the solver-specific representation, rather
than on the normal value}

Issue: DeltaBlue can operate on whole objects, not just primitives
like floats.  How does that get accommodated?
\tf{The way it currently works in Babelsberg/JS is that when a
  DeltaBlue constraint is constructed, and an ``=='' is encountered
  with two variables of the same type on either side, I don't descend
  further, but instead assert the equality on those variables. So for
  example for two points the constraint expression does not desugar
  into multiple constraints on the x and y variables.}


\ab{mirror actual babelsberg implementation, except that we
  re-evaluate the always constraints each time}

\section{Additional Issues}

\emph{This section is currently a holding pen for things that aren't
  workd out.  These should eventually be worked out, and the text
  moved to another part of the memo; or deleted.}

Actual Babelsberg/R code often involves creating new objects in
constraints.  For example, the above constraints on the center of a
rectangle are more clearly written as

\begin{verbatim}
always r.center = Point.new(100,50);
\end{verbatim}

We could view this as sugar for:

\begin{verbatim}
temp := Point.new(100,50);
always r.center = temp;
\end{verbatim}

Similarly, expressions involving point arithmetic create new points on
the fly:

\begin{verbatim}
always r.center = p1+p2;
\end{verbatim}

The interpreter converts this on the fly to:

\begin{verbatim}
always r.center.x = p1.x+p2.x;
always r.center.y = p1.y+p2.y;
\end{verbatim}

An alternative: there are some objects that are like immutable records
(e.g. points), so Babelsberg/Records semantics works.  ID of these objects
not significant.  Do we care about sharing?

How should we handle arrays?  There are a number of odd cases that
arise that are analogous to some of the problems that arise without
structural compatability checks.  (Also see Appendix
\ref{sec:structural-compatability-check-alternatives}.)  Consider:
\begin{verbatim}
Array a;
strong a[10] = 100;
weak a.length = 2;
\end{verbatim}
If this is legal, we want to satisfy the strong constraint in
preference to the weak one.  How big is \verb|a|?  Assuming array
indices start at 1, a length of 10 seems more reasonable than 11 or
100.  But we need a minimality condition to enforce this.

Proposal: the length of an array is part of its structure, and we use
the same structural compatability checks for the existence of array
elements that we do for the existence of fields in a record.  The
above example would then be illegal --- we need to initialize \verb|a|
to be an array of some particular length.  This seems to work and
solve the problem.  The odd thing about it is that the structure of a
record is like its type, and we normally don't consider the length of
an array to be part of its type.  To fix this, we should think of the
structure as being a more stringent notion of type: it includes the
usual properties of a type, plus some others that are usually checked
at runtime.  (As another example, if we are using records to represent
linked lists, the length of the linked list is part of its structure.)  

\tf{I don't think this is very odd -- it makes me think of arrays in
  Smalltalk, where arrays are just objects that are initialized with a
  fixed number of instance variables. the number of instance
  variables, even if they are indexed instead of named, naturally
  ``feels'' like structure to me.}

\ab{Check formal semantics for CLP, cc and tcc languages (which also have
a notion of a constraint store) --- while we are doing something
different in the Babelsberg semantics, a reader who knows that
literature shouldn't come away thinking we are missing something
obvious, e.g. in terms of a better way to do something, falling into a
well-known pitfall, etc.}


\bibliographystyle{plain}
\bibliography{constraints}

\appendix
\newpage

\section{Alternatives to Structural Compatability Checks}
\label{sec:structural-compatability-check-alternatives}

Our model includes structural compatibility checks on constraints,
which allow us to desugar all of the record constraints (recursively)
into constraints on the individual components.  This means that the
solver need not know anything about records.  The alternative of
having the solver know about records was rejected, since it is more
complex and seems to lead to some odd cases, including requiring that
the constraint solver spontaneously generate new records.  For
completeness, we outline that rejected alternative in this subsection.

We first consider records in Babelsberg/Records (immutable records, no UIDs).

If the solver does know about records, we need to extend the
comparators to handle them.  For LPB, there are at least two
possibilities:
\begin{enumerate}
\item An equality constraint between records is either satisfied or
  it's not. For example consider a constraint \verb|p=q| when
  \verb|p={x:0, y:1}| and \verb|q={x:0, y:2}|, versus when
  \verb|q={a:1000}|.  In both cases the constraint is unsatisfied, and
  there is no reason to prefer one of these solutions over the other.

\item An equality constraint between records is unfolded into multiple
  constraints on the fields, each of which is satisfied or not.  In
  the above example, for the first case we have \verb|p.x=q.x|
  (satisfied) and \verb|p.y=q.y| (unsatisfied).  For the second case
  we have \verb|p.x=q.x| (unsatisfied) and \verb|p.y=q.y|
  (unsatisfied), and an additional constraint \verb|p.a=q.a|
  (unsatisfied).  The first solution would be preferred over the
  second under LPB.
\end{enumerate}

We also considered WSB for records.  The error for \verb|{x:5}={x:5}|
or \verb|{x:5}={x:6}| is clear enough (0 and 1 respectively).  But
what is the error for \verb|{x:5}={x:5,y:6}| or \verb|{x:5}={y:6}|?
We need to determine an error for additional or missing fields, and
this error needs to be such that it can be meaningfully compared with
the error for values of fields.  Is the error for
\verb|{x:0}={x:1000000000}| less than for \verb|{x:0}={y:0}|?  We
could come up with a definition, but it's not clear it's that useful.

Another problem with having the solver handle records is that it may
result in less predictability for programmers.  Consider:

\begin{verbatim}
p := {x:0}
q := {y:1}
always p=q
\end{verbatim}

Here there are a number of plausible values for {\tt p} and {\tt q} ---
namely \verb|{x:0}|, \verb|{y:1}|, and \verb|{x:0,y:1}| --- and unclear
which is better.  The nondeterminism here is arguably worse than what
arises with integers and floats and other values, since it affects the
structure of the result rather than just its value.  These issues
probably get even worse once we have heap-allocated records with
aliasing.

Yet another complication is that we would want a minimality condition
on records.  Consider:
\begin{verbatim}
p.x = 5;
\end{verbatim}
We want the solution to be \verb|p={x:5}|, but not e.g.
\verb|p={x:5, y:1000}|.

These issues, along the lack of clear use cases for having a solver
that can solve for records, led us to reject the alternative outlined
in this subsection.

If we move to Babelsberg/UID, an additional odd feature of having the
solver know about records rather than using structural compatability
checks is that the solver may need to spontaneously generate new
records with their own UIDs\@.  Consider:

\begin{verbatim}
a := {x:1};
b := a;
always a.x=1;
always b.x=2;
\end{verbatim}

At first, {\tt a} and {\tt b} refer to the same record, but after the
second {\tt always} constraint, we need to create a new record and
point {\tt b} to it.

For Babelsberg/Objects, the analog of this behavior would likely be
requiring that the constraint solver generate new classes on the fly;
there will also be multiple possible classes.  As before, this seems
overly complex and unpredictable.

\end{document}
